% Section 2.3, from lectures/L02/README.md: the objectives, the evaluation questions, and the next
% lecture.

\section{Review}
\label{c2:sec:review}

\needspace{6\baselineskip}
\subsection*{What you should be able to do now}
After this chapter, you should be able to:
\begin{itemize}
  \item \textbf{Draw a \code{tx} waveform} for a given byte, baud rate and frame format, and mark
    the bit boundaries on it.
  \item \textbf{Explain why the transmitter needs a baud tick at \code{16 x baud}} rather than at
    the baud rate itself, given that the receiver of \chapref{3} oversamples.
  \item \textbf{Implement a transmitter state machine} that loads, shifts and frames, with correct
    \code{busy} and \code{done} timing, verified against a provided testbench.
\end{itemize}

\needspace{6\baselineskip}
\subsection*{Questions to test yourself}
You should be able to explain:
\begin{enumerate}
  \item Why is the line idle-high, and what does a receiver detect to know a frame has started?
  \item A byte is sent at 115200~baud with the 50\,MHz clock: how many system cycles long is one
    bit, and what \code{BAUD_DIV} produces it?
  \item If \code{busy} rises one cycle late, what does the TX feeder of \chapref{5} do on that
    cycle, and what happens to the byte it acts on?
\end{enumerate}

\needspace{6\baselineskip}
\subsection*{Looking ahead}
\chapref{3} takes the hard direction: recovering bytes from an asynchronous line, with a two-flop
synchronizer on \code{rx}, 16x oversampling, mid-bit sampling and framing-error detection.
